% =====================================================================
%  IT4868E Web Mining - Capstone Project Report
%  Compile: latexmk -xelatex report.tex   (or: xelatex, bibtex, xelatex, xelatex)
% =====================================================================
\documentclass[11pt,a4paper]{article}

\usepackage{fontspec} % Latin Modern (supports Vietnamese diacritics)
\usepackage[a4paper,margin=2.4cm]{geometry}
\usepackage{amsmath,amssymb,amsthm}
\usepackage{booktabs,tabularx,multirow,array}
\usepackage[table]{xcolor}
\usepackage{graphicx}
\usepackage{tikz}
\usetikzlibrary{arrows.meta,positioning,shapes.geometric,fit,calc}
\usepackage{pgfplots}
\pgfplotsset{compat=1.17}
\usepackage[ruled,vlined,linesnumbered]{algorithm2e}
\usepackage{enumitem}
\usepackage{caption}
\usepackage[skins,breakable]{tcolorbox}
\usepackage{fancyhdr}
\usepackage[numbers,sort&compress]{natbib}
\usepackage{microtype}
\usepackage[colorlinks=true,linkcolor=blue!50!black,citecolor=green!40!black,urlcolor=blue!60!black]{hyperref}
\usepackage[capitalise,noabbrev]{cleveref}

% draft helpers
\definecolor{todoframe}{RGB}{214,120,0}
\newcommand{\todo}[1]{{\color{todoframe}\textbf{[TODO:} #1\textbf{]}}}
\newcommand{\tbd}{{\color{todoframe}\textbf{TBD}}}
\newtcolorbox{pending}[1][]{colback=orange!6,colframe=todoframe,coltitle=white,boxrule=0.6pt,
  arc=2pt,left=6pt,right=6pt,top=4pt,bottom=4pt,title={#1},fonttitle=\bfseries\small,breakable}

% math 
\newtheorem{definition}{Definition}
\newtheorem{proposition}{Proposition}
\newtheorem{lemma}{Lemma}
\DeclareMathOperator{\regd}{reg}
\DeclareMathOperator{\dom}{dom}
\DeclareMathOperator{\FP}{FP}
\newcommand{\R}{\mathbb{R}}
\newcommand{\one}{\mathbf{1}}

% header
\pagestyle{fancy}
\fancyhf{}
\lhead{\small IT4868E Web Mining --- Capstone Report}
\rhead{\small\color{todoframe}Draft v0.1 --- 28 Oct 2026}
\cfoot{\thepage}
\setlength{\headheight}{14pt}

\title{\textbf{Mining the Third-Party Web in Vietnam:\\
Tracker Ecosystem Discovery via Inclusion-Graph Analysis}\\[6pt]
\large IT4868E Web Mining --- Capstone Project Report (Method type)\\
\normalsize Topic area: Link Analysis}
\author{Baumelou Paul Marie Hadrien 20260145S \and Nguyen Dao Minh Huy 20260104S \\[4pt]
\small Hanoi University of Science and Technology (HUST), semester 2026-1}

\begin{document}
\maketitle

%abstract
\begin{abstract}
Visiting a web page silently triggers requests to dozens of third parties, many of which load
further third parties, forming \emph{inclusion chains} through which user data flows.
Public tracker knowledge bases are compiled from Western vantage points: DuckDuckGo Tracker Radar,
the reference used here, covers nine regions, none of them in Asia. We study the third-party
ecosystem of the Vietnamese web as a \emph{link-analysis} problem. We crawl the 7\,665 most
visited websites in Vietnam (Chrome UX Report country list) with an instrumented headless browser
that records, for every request, the script or document that initiated it. From these logs we
build a domain-level inclusion graph (who loads whom, weighted by the number of sites) and a
co-occurrence graph, and apply the methods of the course---PageRank, HITS, community detection,
node2vec and graph convolutional networks---to (T1) classify unlisted third-party domains as
trackers, (T2) recover their owning organisations, and (T3) mine tracker co-deployment rules.
\end{abstract}

% INTRODUCTION

\section{Introduction}\label{sec:intro}

A modern web page is an aggregate of resources served by many organisations. A news site may
embed a tag manager, which injects an analytics library, which in turn sends identifiers to an
advertising exchange, which redirects to further partners. The site owner explicitly includes
only the first link of this chain; every subsequent party is trusted \emph{implicitly}
\cite{ikram2019chain}. Measuring this ecosystem matters for privacy (who receives user data),
security (each inclusion is an attack surface \cite{kumar2017tangled,arshad2016include}) and
regulation.

Large-scale measurements \cite{englehardt2016online} and curated knowledge bases
(Tracker Radar \cite{trackerradar}, EasyPrivacy \cite{easyprivacy}) describe the Western
web well. Tracker Radar, however, is crawled from nine regions (AU, CA, CH, DE, FR, GB, NL, NO,
US), none in Asia (\cref{sec:data}). Domains that are specific to the Vietnamese advertising
market can therefore be absent from these lists.

\paragraph{Why link analysis.} The syllabus restricts capstone topics to Link Analysis, Opinion
Mining and Recommendation Systems (Information Retrieval excluded). Our object of study is a
\emph{graph}: nodes are third-party domains, and a directed edge $u\to v$ states that $u$ caused
the browser to fetch $v$. This graph plays the role that hyperlinks play in classical link
analysis: a domain loaded by many central loaders is ``important'' in the same sense as a page
cited by authoritative pages. Request URLs are used only as \emph{node attributes}; the project
contains no query processing or document retrieval.

\paragraph{Research questions.}
\begin{description}[leftmargin=1.2cm,labelwidth=1cm,font=\normalfont\bfseries]
  \item[RQ1] How deep and how concentrated is third-party inclusion on the Vietnamese web? Does a
  graph-based ranking (PageRank, HITS) differ from simple prevalence counting?
  \item[RQ2] Can the inclusion and co-occurrence structure, combined with request features,
  identify trackers that are absent from Western lists (T1)?
  \item[RQ3] Does community structure recover the organisations that own the domains (T2)?
  \item[RQ4] Which trackers are systematically co-deployed (T3), and how does prevalence differ
  from the global reference?
\end{description}